\section{Erd\H{o}s Problem 271 --- Round 2}

\subsection*{1) ROUND-2 OBJECTIVE}
Path \textbf{(C)}: the general problem is still open, so the most promising strict advance is to correct the Round~1 computational slip, then extend the fully explicit theory from $A(1)$ to the next nontrivial seed $A(2)$ by a complete proof. This adds a genuinely new explicit family beyond Round~1.

\subsection*{2) Round-1 FOUNDATION USED}
From Round~1 I use only the following established facts.

\begin{itemize}
\item $A(1)$ is exactly the set $S$ of nonnegative integers whose base-$3$ expansion uses only the digits $0$ and $1$.
\item The key representation lemma for $S$: for every integer $u\ge 0$ there exist $a,b\in S$ with
\[
u=2b-a,
\]
and if $u\notin S$ then one may choose $b<u$.
\item The growth of $A(1)$ is $a_k=\Theta(k^{\log_2 3})$.
\end{itemize}

\subsection*{3) NEW INSIGHT / TOOL (ROUND-2)}
The new object is
\[
T:=3S\cup(3S+2).
\]
Equivalently, $T$ is the set of nonnegative integers whose ternary expansion has digits $0$ or $1$ in every position except the units digit, which is allowed to be $0$ or $2$.

The Round-2 claim is that
\[
A(2)=T.
\]
This yields an explicit binary-to-ternary formula and an exact two-line recursion for the sequence generated by $\{0,2\}$.

\subsection*{4) ATTACK PLAN (ROUND-2)}
There were two Round~1 gaps relevant here.

\begin{itemize}
\item The small illustrative sample for $A(4)$ in Round~1 was computational only, and needed verification.
\item No seed beyond $n=1$ had been treated explicitly.
\end{itemize}

The plan is therefore:

\begin{itemize}
\item first correct the $A(4)$ sample by direct greedy checking;
\item then prove that $T=3S\cup(3S+2)$ is 3-free;
\item finally show that every integer $m>2$ outside $T$ is rejected by the greedy algorithm because it completes a 3-term progression with earlier elements of $T$.
\end{itemize}

\subsection*{5) WORK (ROUND-2)}

\paragraph{Round-2 correction of the Round-1 sample.}
The Round~1 sample for $A(4)$ was incorrect. The correct initial segment is
\[
A(4):\ 0,4,5,7,11,12,16,23,26,31,33,37,38,44,49,56,\dots
\]
A quick check already shows why $6$ could not have appeared: after choosing $0,4,5$, the candidate $6$ completes the arithmetic progression $4,5,6$ and must be rejected.

This correction affects only the heuristic illustration from Round~1; none of the actual theorems from Round~1 depended on the erroneous list.

\paragraph{Theorem 5.1.}
Let $S=A(1)$, and define
\[
T:=3S\cup(3S+2).
\]
Then the Stanley sequence generated by $\{0,2\}$ is exactly $T$.

\emph{Proof.}
We first show that $T$ is 3-free.

Suppose $x<y<z$ lie in $T$ and satisfy $x+z=2y$.
Reduce modulo $3$.

\begin{itemize}
\item If $x\equiv z\equiv 0\pmod 3$, then also $y\equiv 0\pmod 3$, so
\[
x=3x',\quad y=3y',\quad z=3z'
\]
with $x',y',z'\in S$, and then $x'+z'=2y'$, contradicting the 3-freeness of $S$.

\item If $x\equiv z\equiv 2\pmod 3$, then also $y\equiv 2\pmod 3$, so
\[
x=3x'+2,\quad y=3y'+2,\quad z=3z'+2
\]
with $x',y',z'\in S$, and again $x'+z'=2y'$, contradiction.

\item If $x\equiv 0\pmod 3$ and $z\equiv 2\pmod 3$ (or vice versa), then
\[
x+z\equiv 2\pmod 3,
\]
so $2y\equiv 2\pmod 3$, hence $y\equiv 1\pmod 3$, impossible because every element of $T$ is congruent to $0$ or $2$ modulo $3$.
\end{itemize}

Thus $T$ contains no nontrivial 3-term arithmetic progression.

Now let $m>2$ be an integer with $m\notin T$. We show that $m$ is rejected by the greedy algorithm for seed $\{0,2\}$.
There are three residue classes modulo $3$.

\smallskip
\noindent
\emph{Case 1: $m\equiv 0\pmod 3$.}
Write $m=3u$. Since $m\notin T$, we have $u\notin S$.
By the Round~1 representation lemma for $S$, there exist $a,b\in S$ with
\[
u=2b-a,
\qquad b<u.
\]
Then
\[
m=3u=2(3b)-3a.
\]
Both $3a$ and $3b$ lie in $T$, and $3b<m$ because $b<u$. Hence $m$ completes a 3-term progression with earlier elements of $T$.

\smallskip
\noindent
\emph{Case 2: $m\equiv 2\pmod 3$.}
Write $m=3u+2$. Since $m\notin T$, again $u\notin S$.
Choose $a,b\in S$ with $u=2b-a$ and $b<u$.
Then
\[
m=3u+2=2(3b+2)-(3a+2).
\]
Now $3a+2,3b+2\in T$, and $3b+2<m$ because $b<u$. So $m$ is rejected.

\smallskip
\noindent
\emph{Case 3: $m\equiv 1\pmod 3$.}
Write $m=3u+1$. Since $m>2$, we have $u\ge 1$.
Apply the Round~1 representation lemma to $u-1$: choose $a,b\in S$ with
\[
u-1=2b-a,
\qquad b\le u-1.
\]
Then
\[
m=3u+1=3(u-1)+4=2(3b+2)-3a.
\]
Again $3a\in T$ and $3b+2\in T$, while
\[
3b+2\le 3(u-1)+2=3u-1<m.
\]
So $m$ is rejected.

Therefore every integer $m>2$ outside $T$ is rejected, while every integer in $T$ is accepted because $T$ is 3-free. Since the seed is $\{0,2\}$, the greedy sequence is exactly $T$. 

\paragraph{Corollary 5.2 (explicit formula for $A(2)$).}
Let
\[
A(2)=\{b_0<b_1<b_2<\cdots\}.
\]
If
\[
k=\sum_{i\ge 0}\varepsilon_i 2^i
\qquad (\varepsilon_i\in\{0,1\})
\]
is the binary expansion of $k$, then
\[
b_k=2\varepsilon_0+\sum_{i\ge 1}\varepsilon_i 3^i.
\]
Equivalently, if $A(1)=\{a_0<a_1<a_2<\cdots\}$, then
\[
b_{2m}=3a_m,
\qquad
b_{2m+1}=3a_m+2.
\]

\emph{Proof.}
The set $T=3S\cup(3S+2)$ is the disjoint union of two increasing translates of $S$, and for every $s\in S$ we have
\[
3s<3s+2<3s'
\qquad\text{whenever } s'\ge s+1.
\]
Hence the increasing enumeration of $T$ interleaves as
\[
3a_0,\ 3a_0+2,\ 3a_1,\ 3a_1+2,\ 3a_2,\ 3a_2+2,\dots,
\]
which is exactly the stated recursion. Rewriting this in terms of the binary digits of $k$ gives the explicit formula. 

\paragraph{Corollary 5.3 (growth of $A(2)$).}
The sequence $A(2)$ satisfies
\[
b_k=\Theta(k^{\log_2 3}).
\]

\emph{Proof.}
From Corollary~5.2,
\[
b_{2m}=3a_m,
\qquad
b_{2m+1}=3a_m+2.
\]
Since Round~1 proved $a_m=\Theta(m^{\log_2 3})$, the same asymptotic follows immediately for $b_k$. 

\subsection*{6) ADVERSARIAL VERIFICATION}
\begin{itemize}
\item The correction of the $A(4)$ sample is genuine: the earlier list was wrong, but only in a heuristic example, not in any proved statement.
\item In the proof of Theorem~5.1, the only nontrivial imported fact is the Round~1 representation lemma for $S=A(1)$. I used it exactly in the two forms already proved there: for $u\notin S$, one may choose $u=2b-a$ with $b<u$; and for arbitrary $u\ge 0$, one may choose $u=2b-a$ with $b\le u$.
\item The residue-$1$ case is the delicate one. The formula
\[
3u+1=2(3b+2)-3a
\]
requires a representation of $u-1$, not of $u$. This was checked explicitly; the inequality $3b+2<m$ follows from $b\le u-1$.
\item The proof is self-contained once Round~1 is accepted.
\end{itemize}

\subsection*{7) FINAL (EXACTLY ONE)}
\textbf{UNRESOLVED (BUT STRICTLY ADVANCED)}

Round~2 corrects the faulty $A(4)$ sample and gives a full explicit determination of $A(2)$, including the exact recursion
\[
b_{2m}=3a_m,
\qquad
b_{2m+1}=3a_m+2,
\]
and the growth law $A(2)=\Theta(k^{\log_2 3})$. The general problem for arbitrary $n$ remains open.

\subsection*{8) COMPLETION ESTIMATE}
COMPLETION: 65\%

\subsection*{9) REFERENCES}
\begin{itemize}
\item A. M. Odlyzko and R. P. Stanley, \emph{Some Curious Sequences Constructed with the Greedy Algorithm}, January 1978. Checked in this round for the classical regular/irregular context and the corrected initial examples.
\item T. F. Bloom, \emph{Erd\H{o}s Problem \#271}, checked 2026-03-07.
\end{itemize}

